\documentclass[11pt,a4paper]{article}
\usepackage{amsmath,amssymb}
\usepackage{listings}
\usepackage{hyperref}
\usepackage[margin=1in]{geometry}
\usepackage{fancyhdr}

\pagestyle{fancy}
\fancyhf{}
\fancyhead[L]{Module 13: Debugging with GDB}
\fancyhead[R]{C Programming Course}
\fancyfoot[C]{\thepage}

\lstset{
  language=C,
  basicstyle=\ttfamily\small,
  keywordstyle=\bfseries,
  commentstyle=\itshape,
  numbers=left,
  numberstyle=\tiny,
  frame=single,
  breaklines=true,
  tabsize=4
}

\title{Module 13: Debugging with GDB}
\author{C Programming Course}
\date{}

\begin{document}
\maketitle
\thispagestyle{fancy}

\section{Introduction to GDB}
GDB (GNU Debugger) allows you to step through code, inspect variables, and find
bugs at runtime. Compile with \texttt{-g} to include debug symbols:
\begin{verbatim}
gcc -g -o program program.c
gdb ./program
\end{verbatim}

\section{Setting Breakpoints}
\begin{verbatim}
(gdb) break main
(gdb) break myfile.c:42
(gdb) break my_function
(gdb) info breakpoints
(gdb) delete 1
\end{verbatim}

Conditional breakpoints:
\begin{verbatim}
(gdb) break myfile.c:30 if i == 5
\end{verbatim}

\section{Stepping Through Code}
\begin{itemize}
  \item \texttt{run} -- start the program.
  \item \texttt{next} (n) -- execute the next line, stepping over function calls.
  \item \texttt{step} (s) -- step into function calls.
  \item \texttt{continue} (c) -- resume until the next breakpoint.
  \item \texttt{finish} -- run until the current function returns.
\end{itemize}

\section{Examining Variables and Memory}
\begin{verbatim}
(gdb) print x
(gdb) print *array@10
(gdb) x/4xw &variable
(gdb) display x
(gdb) info locals
\end{verbatim}

\section{Advanced GDB Techniques}
\subsection{Watchpoints}
Break when a variable changes:
\begin{verbatim}
(gdb) watch my_var
\end{verbatim}

\subsection{Backtrace}
Show the call stack after a crash:
\begin{verbatim}
(gdb) backtrace
\end{verbatim}

\section{Debugging Memory Issues}
Example: detecting an out-of-bounds access.

\begin{lstlisting}
#include <stdio.h>
#include <stdlib.h>

int main(void) {
    int *arr = malloc(5 * sizeof(int));
    if (arr == NULL) return 1;
    for (int i = 0; i < 5; i++) {
        arr[i] = i;
    }
    /* Use GDB to inspect arr contents */
    printf("arr[2] = %d\n", arr[2]);
    free(arr);
    return 0;
}
\end{lstlisting}

Compile with \texttt{-g -fsanitize=address} to catch memory errors automatically.

\end{document}
